\section{Coefficient-selected rectangular means for arbitrary scalar profiles}
\label{sec:balanced223}
This section supplies the finite component-gradient construction asked for
in Section~27. It uses the original positive-start scalar coefficient and
retains the edges crossing rectangular observation pieces inside a ring.
It does not identify a variable-weight curl/divergence form with that gradient
form. The balanced-ring, BV and tree mechanisms are classical;
historical originality of this finite specialization is unestablished.

\subsection{Statement and actual observations}
Let integers $D\ge3$, $n\ge1$, and $c>0$ be fixed, and put
$X=\{0,\ldots,n-1\}^D$. For any finite scalar profile $a(x)>0$ define
\[
 P=D/2,\quad b=2^D,\quad J=2D,\qquad
 C_{\rm ext}=1+J\{2J(1+12J)+1\},\qquad
 G_D=2^{2D+10}D^2C_{\rm ext}^{\,2}.
\]
Write $S(A)=\sum_{x\in A}a(x)^{-P}$ and
$t=(cG_D)^{-P}/b$. Every internal edge below keeps its original positive
start, and all sums and norms initially use counting measure.
\begin{theorem}\label{thm:balanced223}
There is a partition $\mathcal R$ of $X$ into nonempty axis-aligned integer
rectangles, determined by $a$ and fixed geometric conventions alone, such that
\[
 \#\mathcal R\le2D\{1+(2b+3)S(X)/t\},
\]
and, simultaneously for every real field $f$,
\begin{equation}\label{eq:balanced223-final}
 \sum_{i,x:\ x,x+e_i\in X}a(x)|f(x+e_i)-f(x)|^2
 +c\sum_{B\in\mathcal R}\frac1{|B|}\Bigl(\sum_{x\in B}f(x)\Bigr)^2
 \ge c\sum_{x\in X}|f(x)|^2.
\end{equation}
The original periodic energy also satisfies this inequality, by adding
its nonnegative wrap terms. For $h=D(N_c^2-1)$ real components the same
partition has output dimension $h\#\mathcal R$ and penalty $\Phi(y)=c\|y\|^2$.
\end{theorem}
\begin{corollary}\label{cor:balanced223-moment}
If $\mu\ge1$ and
$n^{-D}\sum_x(\mu/a(x))^{D/2}\le M$, the actual normalized rectangular
means satisfy the rank order in \eqref{ap:eq:rank}, with $K=c$ and safe constant
\[
 C_{\rm rect}=2Dh\{1+(2b+3)b(cG_D)^P2^PM\}.
\]
The exact arithmetic mean and envelope $1\le a\le2\mu$ are unnecessary
for this sufficiency. The inverse-moment hypothesis is still assumed;
it is not deduced from an arithmetic mean.
\end{corollary}
\begin{proof}
$S(X)\le Mn^D\mu^{-P}\le2^PM\max\{1,n^D/(1+\mu)^P\}$, since
$\mu\ge1$. Insert $1/t=b(cG_D)^P$ into the theorem. Apply its scalar
inequality to each of the $h$ real components. Dividing energy and squared
norms by $V=n^D$ gives observation rows
$V^{-1/2}|B|^{-1/2}\sum_B f=\sqrt{|B|/V}\,\operatorname{mean}_B f$.
The all-field inequality restricts to divergence-zero fields, and linear
means with the displayed quadratic penalty supply the original local data.
\end{proof}

\subsection{Balanced integer nodes and actual collars}
Split every coordinate interval with at least two sites into its floor/ceiling
halves; a singleton remains unchanged. Use coordinate order $1,\ldots,D$,
lower before upper, and lexicographic order on children. Each node has at
most $b$ children, and all side lengths differ by at most one. At global
depth $k$ the nonempty intervals have lengths $h_k$ or $h_k+1$, until they
are singletons, where $h_k=\max\{1,\lfloor n/2^k\rfloor\}$.

Consider a proper descendant $I$ of an ancestor $Q$, at global depth $k$,
and set $h_0=h_k$. Every side of $I$ belongs to $[h_0,2h_0]$.
Every positive clearance from $I$ to the boundary of $Q$ is a union of
depth-$k$ intervals, so is at least $h_0$. This uses the global $h_0$,
not the smallest side of this particular $I$: at odd sizes that substitution
can be false. If every side of $Q$ is at least two, a proper descendant
is shorter in every direction, so each axis has at least one positive
clearance. If the smallest side of $Q$ is one, every side is one or two;
each proper descendant is a corner singleton. That small case is treated
separately below.

For the large case put $\Omega=Q\setminus I$. Expand $I$ by exactly $h_0$
on every face with positive clearance, leaving a boundary-touching face
unchanged, and call the resulting rectangle $B$. Its side lengths $\ell_i$
lie in $[2h_0,4h_0]$. The collar $A=B\setminus I$ is covered by the actual
positive-clearance face slabs: normal width $h_0$ and all other sides those
of $B$. There are between $D$ and $J=2D$ slabs, each contained in $\Omega$.
For slabs $S_i,S_j$ with distinct normals, direct product counting gives
\begin{equation}\label{eq:balanced223-overlap}
 \frac{|S_i|}{|S_i\cap S_j|}=\frac{\ell_j}{h_0}\le4,
 \qquad \frac{|S_i|}{|S_j|}=\frac{\ell_j}{\ell_i}\le2.
\end{equation}
Choose the first slab as root; attach every distinct-normal slab directly
to it, and its opposite slab, if present, through the first slab with a
different normal. This tree has depth at most two.

\subsection{An unweighted BV extension, with every hole-crossing edge}
For a finite product rectangle $S$ of longest side at most $4h_0$, write
$V_1(S)=\sum_S|f-\operatorname{mean}_S f|$ and let $\mathrm{TV}_S(f)$
be the sum of $|df|$ over its actual internal edges. Then
\begin{equation}\label{eq:balanced223-L1}
 V_1(S)\le2h_0\mathrm{TV}_S(f).
\end{equation}
Indeed on an interval with $m$ sites, averaging pair paths bounds the
coefficient of its $r$th edge by $2r(m-r)/m\le m/2$.
Commuting coordinate-average projections contract $L^1$; telescoping their
product and applying this interval inequality fibrewise proves
\eqref{eq:balanced223-L1}, without an additional factor $D$.

Let $m_i$ be the mean on $S_i$ and let $c_0$ be the root mean.
Averaging on the actual intersection and using
\eqref{eq:balanced223-overlap} gives
$|S_i||m_i-m_j|\le4\{V_1(S_i)+V_1(S_j)\}$.
For a depth-two path $i,p,r$, the second edge costs the volume ratio
$|S_i|/|S_p|\le2$, so
\[
 |S_i||m_i-m_r|
 \le4V_1(S_i)+12V_1(S_p)+8V_1(S_r)
 \le12\sum_j V_1(S_j).
\]
Direct paths satisfy the same safe upper bound. Covering $A$ and counting
each actual edge at most $J$ times yields, with $P_D=2J(1+12J)$,
\begin{equation}\label{eq:balanced223-collar}
 \sum_A|f-c_0|\le(1+12J)\sum_jV_1(S_j)
 \le P_Dh_0\mathrm{TV}_A(f).
\end{equation}

At each positive-clearance hole face take the adjacent outside layer
of width $h_0$, with tangential sides those of $I$. It is contained in
$A$. On a normal fibre its nearest value satisfies
\[
 |f(\mathrm{first})-c_0|
 \le h_0^{-1}\sum_{\mathrm{fibre}}|f-c_0|
       +\sum_{\mathrm{fibre}}|df|.
\]
Sum over all actual faces and use \eqref{eq:balanced223-collar}.
The sum of hole-crossing jumps to the constant $c_0$ is at most
$J(P_D+1)\mathrm{TV}_A(f)$. Fill $I$ by $c_0$ and leave $f$ unchanged
on $\Omega$, obtaining $\widetilde f$ on $Q$. Internal hole differences
vanish; the added variation consists precisely of these crossing faces.
Therefore
\begin{equation}\label{eq:balanced223-BV}
 \mathrm{TV}_Q(\widetilde f)
 \le C_{\rm ext}\mathrm{TV}_\Omega(f).
\end{equation}
Boundary-touching faces create no exterior edge; at $h_0=1$ the fibre
gradient sum is zero. No coefficient is extended, reflected, averaged or
reassigned. In particular \eqref{eq:balanced223-BV} makes no weighted
quadratic trace-extension claim of the kind excluded earlier.

\subsection{Actual critical Sobolev and weighted gap on a ring}
Here are the finite details needed to pass from BV to the original energy.
For $T\subset\{0,\ldots,m-1\}^D$ with $|T|\le m^D/2$, the occupied-fibre
argument and the discrete projection inequality give
\[
 |\partial T|\ge\sum_i|\pi_iT|-D|T|/m
 \ge D|T|^{(D-1)/D}-D|T|/m
 \ge c_{\rm iso}|T|^{(D-1)/D},\quad
 c_{\rm iso}=D(1-2^{-1/D})\ge\tfrac12.
\]
Each occupied nonfull fibre supplies an edge, while full fibres number at
most $|T|/m$. The projection inequality
$|T|^{D-1}\le\prod_i|\pi_iT|$ is the finite Loomis--Whitney inequality;
it follows from iterated H\"older, or the finite projection argument used
in the earlier Sobolev section. A singleton cube has only zero variance.
Median parts, finite coarea and Minkowski consequently imply
\[
 \inf_z\|u-z\|_{D/(D-1)}\le(2/c_{\rm iso})\mathrm{TV}(u)
\]
on the cube. On a rectangle of aspect at most two, fold a cube of its
longest side onto it: in each coordinate first visit the interval in
order, repeat its final endpoint and reflect until the prescribed length.
Every vertex and edge has at most two preimages per coordinate; the
initial product block maps identically onto the rectangle. Restriction
and the variation multiplicity at most $2^D$ give the safe constant
$C_{\rm full}=2^{D+1}/c_{\rm iso}$ for that rectangle, still unweighted.

Apply this bound and \eqref{eq:balanced223-BV} to an indicator on $\Omega$.
For every real $z$,
$\|1_T-z\|_{r,\Omega}\ge\frac12\min\{|T|,|\Omega\setminus T|\}^{1/r}$,
where $r=D/(D-1)$. Thus $\Omega$ has relative edge isoperimetric constant
\[
 c_\Omega=c_{\rm iso}/(2^{D+2}C_{\rm ext}).
\]
Coarea gives $\|g\|_r\le c_\Omega^{-1}\mathrm{TV}_\Omega(g)$ for
nonnegative $g$ supported on at most half the vertices. With a median $m_f$
of $f$, take $w=(f-m_f)_+$ or $(m_f-f)_+$ and put
\[
 \alpha=2(D-1)/D,\qquad q=2D/(D+2),\qquad q'=2D/(D-2).
\]
Then $\alpha r=2$ and $(\alpha-1)q'=2$. The power-difference inequality,
H\"older and degree at most $2D$ yield
\[
 \mathrm{TV}(w^\alpha)
 \le2\alpha D^{1/q'}\|dw\|_q\|w\|_2^{\alpha-1}.
\]
In fact the $q'$th power of the endpoint factor sums to at most
$2^{q'}D\sum_xw(x)^2$. Divide only when $w\ne0$. The positive and
negative truncated $q$-edge energies sum to at most the original one;
since $2/q\ge1$, their squared bounds combine to give
\begin{equation}\label{eq:balanced223-q}
 \sum_\Omega|f-\operatorname{mean}_\Omega f|^2
 \le K_\Omega^2\Bigl(\sum_{\Omega\text{ edges}}|df|^q\Bigr)^{2/q},
 \qquad K_\Omega=2\alpha D^{1/q'}/c_\Omega.
\end{equation}
Only now use weights, on the original retained positive-start edges:
\[
 \Bigl(\sum |df|^q\Bigr)^{2/q}
 \le\Bigl(\sum_{\Omega\text{ edges}}a(\mathrm{start})^{-D/2}\Bigr)^{2/D}
       \E_a(f;\Omega)
 \le\{DS(\Omega)\}^{2/D}\E_a(f;\Omega).
\]
Here $q/(2-q)=D/2$, and at most $D$ retained edges start at one vertex.
Since $\alpha\le2$, $c_{\rm iso}\ge1/2$ and $2/q'+2/D=1$,
$K_\Omega^2D^{2/D}\le2^{2D+10}DC_{\rm ext}^{\,2}\le G_D$.
Full balanced rectangles obey the same safe bound without removing a hole.

For the small nodes with a singleton direction, $Q$, or $Q$ minus a corner,
is a connected graph with at most $b$ vertices whenever nonempty. Pair paths
of length at most $b-1$ give variance at most $b^2\E_1$.
Each internal edge satisfies
$a(\mathrm{start})^{-D/2}\le S(\Omega)$, so
$\E_1\le S(\Omega)^{2/D}\E_a$. As $G_D\ge b^2$, this proves in every case
\begin{equation}\label{eq:balanced223-gap}
 \sum_\Omega|f-\operatorname{mean}_\Omega f|^2
 \le G_DS(\Omega)^{2/D}\E_a(f;\Omega).
\end{equation}
A singleton has zero variance. Empty groups are never output.

\subsection{Coefficient-only heavy-chain compression and complete count}
Give each vertex mass $a(x)^{-P}$ and call a tree node heavy if its mass is
strictly greater than $t$. If the root is light, output it. Otherwise let
$L$ be the number of terminal heavy nodes and $B_h$ the number with at
least two heavy children. Terminal heavy nodes are disjoint, so
$L\le S(X)/t$. Finite tree counting gives $B_h\le L-1$.
After contracting unary paths, there are at most $B_h+L-1$ skeleton edges
and possibly a root-leading chain; hence at most $B_h+L\le2L-1$ maximal
unary chains. The bound is on chains, not their unrestricted total length.

At a unary node $Q\to Q'$ the increment $Q\setminus Q'$ is the union of
at most $b-1$ light children, of mass at most $(b-1)t$. Traverse a maximal
chain from ancestor toward descendant, closing a group when the cumulative
mass first reaches $t$. A closed group is a balanced ancestor ring
$Q_i\setminus Q_j$ of mass in $[t,bt)$; a residual is output only if nonempty,
and has mass less than $t$. Across all chains the disjoint closed groups
number at most $S(X)/t$, while residuals number at most $2L-1$.
At branching or terminal heavy internal nodes output each light child;
at a heavy singleton terminal output the singleton. Such cubes/singletons
number at most $b(B_h+L)\le b(2L-1)$.

These groups partition $X$: a unary ring excludes the endpoint heavy
subtree, which is then handled at the endpoint; branching processing
recurses into every heavy child and outputs all remaining children.
Thus their total number is at most
\[
 S(X)/t+(2L-1)+b(2L-1)\le1+(2b+3)S(X)/t.
\]
Each nonsingleton group has mass at most
$bt=(cG_D)^{-D/2}$, so \eqref{eq:balanced223-gap} gives residual gap at least
$c$. Heavy singleton mass may exceed this threshold, but its variance is
zero and its mean captures its entire norm.

\subsection{Rectangular rows while retaining intra-ring interfaces}
Partition $Q\setminus I$ by the first coordinate outside $I$: preceding
coordinates belong to $I$, the current coordinate is below or above $I$,
and later coordinates fill $Q$. This gives at most $2D$ nonempty integer
rectangles. Cubes and singletons require one rectangle. With
$q_A(f)=|A|^{-1/2}\sum_A f$, Cauchy gives
$q_\Omega(f)^2\le\sum_{B\subset\Omega}q_B(f)^2$.
The actual regional gap therefore implies
\[
 \E_a(f;\Omega)+c\sum_{B\subset\Omega}q_B(f)^2\ge c\|f\|_\Omega^2.
\]
Crucially $\E_a(f;\Omega)$ retains all edges between rectangular pieces
of this same ring; no separate gap is asserted for each piece.
Distinct regions are disjoint, so summing discards only nonnegative edges
between different regions. The full energy dominates their internal sum.
The count of rectangles is at most $2D$ times the region count, proving
Theorem~\ref{thm:balanced223}. This directly avoids the internal-cell
interface discard excluded by the earlier canonical example.

\subsection{Attribution, audit and executable evidence}
Brudnyi~\cite{brudnyi221}, Sections~1--2 and Appendices~I--II, and its cited
predecessors already develop dyadic-ring approximation and related tree
mechanisms. Cohen--DeVore--Petrushev--Xu~\cite{cdpx223} give a classical
bivariate BV approximation setting, also with dyadic rings; that
function-approximation problem is not asserted to originate here. Remark~I.2 includes boundary-touching inner cubes. Neither
ring decomposition, BV/Sobolev machinery, boundary touching nor heavy-tree
compression is asserted to be a new general principle here. This section
supplies the explicit finite integer coefficient-only original-energy
adapter; its exact historical priority remains unestablished. In particular,
an atomic coefficient mass is handled by terminal singletons, rather than
silently importing continuum absolute continuity.

The original analytical draft is at programme source
\begin{center}\scriptsize
\texttt{c993dff04868d01d37b6ec87e0417da3c91a1330}.
\end{center}
Its raw-byte SHA-256 is
\begin{center}\scriptsize
\texttt{f2cd7941ccc8e9fa30e3a8742d5e814b0e4bcdc83ded16426b48b78c7deac057}.
\end{center}
The full involved AI source audit accepted it with limits and no load-bearing
error; its report, factual shared-context limits and requested minor
clarifications are retained in the source/evidence packet. The explicit
$n\ge1,c>0$, nonempty-residual and fixed-order conventions above address
those clarifications without changing constants or the proof.
Separate own Colab normal and optimized executions at that source had
actual exits zero and identical output: 216 selected geometries,
648 BV checks, 200 coefficient partitions, 71,232 regional-field and
800 global-field inequalities, and three rejected real negative controls.
The counted within-ring cross-piece edges number 1,668. These selected
exact-rational checks do not constitute an all-size executable proof,
formal certification or two independent fresh-clone reproductions.
The runtime was released and separately read back inactive; billed units
are unknown. No local scientific execution was used.

This result closes the stated finite scalar-gradient rectangular
construction analytically. It supplies no covariant/interacting coefficient
producer, variable-weight curl replacement or physical continuum theorem.
PR29 remains at 20/41 with TermSource zero. This changed manuscript requires
its own exact-PDF assessment; no numerical grade transfers from version~12.
